\section{Introduction}
Orange operates large IP/MPLS backbone networks, where \emph{traffic engineering} aims at avoiding congestion and keeping a good quality of service, also when some links are unavailable because of failures or planned maintenance. With \emph{Segment Routing} (SR) a packet carries a list of waypoints: between two consecutive waypoints it follows the usual shortest paths, and when several shortest paths exist the traffic is split evenly at every node (ECMP). Choosing waypoints therefore allows to leave the shortest paths and to balance the load without changing the routers' configuration.

The \emph{$T$-Adaptive Segment Routing} problem (T-ASR) of the Orange challenge is the following. We are given a directed network $G=(V,A)$ with an arc weight $\omega(a)$ and a capacity $c(a)$, a set $D$ of demands $d=(s_d,\tau_d)$ with a traffic volume $\delta(d,t)$ for every time period $t\in T=\{0,\dots,h-1\}$, a maintenance scenario $q(t)\subseteq A$ (arcs that are down at period $t$), a budget $\beta(t)$ limiting the number of changes of the routing between $t-1$ and $t$, and a maximum number of segments $\mathrm{maxSeg}$ per demand. We must choose an SR path for every demand and every period so that the loads of the links are as balanced as possible. In this work we study periods $T=\{0,1\}$: period $0$ is the nominal state (no budget, no failure), and period~$1$ has one link down and a budget.
